\documentclass[12 pt]{abntex2}
\usepackage[utf8]{inputenc} % Codificação de caracteres
% \usepackage[T1]{fontenc} % Seleção de fonte
\usepackage[brazil]{babel} % Idioma do documento
\usepackage{graphicx} %For figures
\usepackage{float} %For figures
\usepackage{amsmath}
\usepackage{amssymb}
\usepackage{rsfso}
\usepackage{mathrsfs}
\usepackage{lmodern}
\usepackage{hyperref}
\usepackage{listings}
\hypersetup{
    colorlinks=true,
    linkcolor=blue,
    filecolor=magenta,
    urlcolor=cyan,
    pdftitle={},
    pdfpagemode=FullScreen,
}


\hbadness=10000 % Suppresses Underfull \hbox warnings
\hfuzz=9999pt  % Suppresses Overfull \hbox warnings

% \usepackage{pgfplots}
% \pgfplotsset{compat=1.18}

% \begin{tikzpicture}
% \begin{axis}[
%     axis lines = middle, % Draws axes through the origin
%     xlabel = $x$,
%     ylabel = {$f(x)$},
%     xmin=-5, xmax=5,
%     ymin=-5, ymax=5,
%     samples=100, % Increase samples for smoother curves
% ]
% \addplot[domain=-5:5] {x^2}; % Plots the function y = x^2
% \end{axis}
% \end{tikzpicture}


\makeatletter
\renewcommand{\thesection}{%
  \ifnum\c@chapter<1 \@arabic\c@section
  \else \thechapter.\@arabic\c@section
  \fi
}
\makeatother

    % \begin{figure}[H]
    %   \centering
    %   \includegraphics[width =0.50\linewidth, angle = 0]{file.png}
    %   \caption{caption about}
    %   \label{} %if necessary
    % \end{figure}
    % Code to add figures

    % \begin{table}[h!]
    %     \centering
    %     \caption{Example Table}
    %     \label{tab:my_table}
    %     \begin{tabular}{|l|c|r|}
    %         \hline
    %         \textbf{Column 1} & \textbf{Column 2} & \textbf{Column 3} \\
    %         \hline
    %         Row 1, Cell 1 & Row 1, Cell 2 & Row 1, Cell 3 \\
    %         Row 2, Cell 1 & Row 2, Cell 2 & Row 2, Cell 3 \\
    %         \hline
    %     \end{tabular}
    % \end{table}
    % Code for adding tables

\autor{Lucas Tasca, Melkior Gabriel Paim Misiak}
\titulo{Avaliação M1 \\ IPC, Threads e Paralelismo}
\orientador{Felipe Viel}
\instituicao{Universidade do Vale do Itajaí - UNIVALI\\ Sistemas Operacionais}
\local{Itajaí - Santa Catarina}
\data{\today}

\begin{document}

\imprimirfolhaderosto

\tableofcontents*
\clearpage

\section{Introdução}

Este projeto apresenta a implementação de um sistema de processamento de
imagens do tipo PGM P5, com o objetivo de consolidar o aprendizado sobre
comunicação entre processos (IPC), threads, concorrência e
paralelismo em sistemas operacionais Unix/Linux.
 
O sistema é dividido em dois executáveis independentes que se comunicam por
uma FIFO (named pipe), o processo \textbf{emissor} (\texttt{img\_sender}) lê uma
imagem PGM do disco e a transmite para a FIFO e o processo \textbf{trabalhador}
(\texttt{img\_worker}) recebe a imagem, distribui o processamento entre
múltiplas \textit{threads} por meio de sincronização de tarefas e grava o resultado em disco.
 
A paralelização é realizada por decomposição de linhas: cada \textit{thread}
processa um subconjunto das linhas da imagem, sem sobreposição, o que garante
correção sem necessidade de sincronização na escrita dos pixels. A fila de
tarefas é protegida por \textit{mutex} e semáforos contadores, implementando
o padrão clássico produtor-consumidor.

\section{Fundamentação Teórica}
 
\subsection{Comunicação entre Processos via FIFO}
 
FIFOs (\textit{named pipes}) são arquivos especiais do sistema de arquivos
Unix que permitem comunicação unidirecional entre processos independentes. A
abertura de um FIFO é bloqueante: o processo que abre para leitura aguarda
até que outro processo abra para escrita, garantindo sincronização implícita
no início da comunicação.
 
\subsection{Threads e Pthreads}
 
\textit{Threads} são unidades de execução que compartilham o mesmo espaço de
endereçamento dentro de um processo. A biblioteca POSIX Threads (Pthreads)
fornece as primitivas necessárias:
 
\begin{itemize}
  \item \texttt{pthread\_create} -- criação de \textit{threads};
  \item \texttt{pthread\_join} -- aguardar o término de uma \textit{thread};
  \item \texttt{pthread\_exit} -- encerramento da \textit{thread} atual.
\end{itemize}
 
\subsection{Semáforos}
 
Semáforos contadores (\texttt{sem\_t}) são mecanismos de sincronização que
controlam o acesso concorrente a recursos compartilhados. Neste projeto é
utilizado um único semáforo \texttt{sem}, inicializado com o número de
\textit{threads} \texttt{g\_nthreads}. Ele regula quantas \textit{threads}
podem acessar simultaneamente a fila de tarefas e executar o processamento,
atuando como guarda de exclusão mútua da seção crítica composta pelo
dequeue e pelo processamento do bloco de linhas.
 
\subsection{Filtros de Processamento de Imagem}
 
Onde $r$ é o valor do pixel de entrada e $s$ é o pixel de saída.

\subsubsection{Filtro Negativo}

Transformação linear que inverte os tons de cinza:

\begin{equation}
  s = \mathtt{maxv} - r
\end{equation}

$\mathtt{maxv}$ é o valor máximo definido no cabeçalho PGM (tipicamente 255).

\subsubsection{Limiarização com Fatiamento}
 
Preserva os pixels dentro do intervalo $[t_1, t_2]$ e zera os demais:
 
\begin{equation}
  s = \begin{cases}
    r & \text{se } t_1 \leq r \leq t_2 \\
    0 & \text{caso contrário}
  \end{cases}
\end{equation}

\section{Arquitetura da Solução}
 
A solução é composta por dois processos independentes e um cabeçalho
compartilhado (\texttt{pgm.h}), conforme descrito a seguir.
 
\subsection{Estruturas de Dados}
 
\textbf{Estruturas principais -- \texttt{pgm.h}}
\begin{lstlisting}[language=C]
typedef struct{
    char type[3];
    uint32_t w, h; 
    uint8_t maxv; 
    uint8_t* data;
} PGM; 

typedef struct{
    uint32_t w, h;
    uint8_t maxv;
    OP_Mode mode; 
    uint8_t t1, t2; 
} Header;

typedef struct{ 
    uint32_t row_start; 
    uint32_t row_end; 
} Task;
\end{lstlisting}
 
\subsection{Fluxo de Dados}
 
\begin{enumerate}
  \item \texttt{img\_sender} lê a imagem PGM do disco com \texttt{read\_pgm};
  \item \texttt{img\_sender} abre o FIFO para escrita (bloqueia até \texttt{img\_worker} abrir para leitura);
  \item \texttt{img\_sender} transmite o \texttt{Header} seguido dos pixels;
  \item \texttt{img\_worker} recebe os dados e aloca \texttt{g\_in};
  \item \texttt{img\_worker} cria $N$ \textit{threads} e enfileira $N$ tarefas (uma por \textit{thread});
  \item Cada \textit{thread} consome uma tarefa, processa as linhas correspondentes em \texttt{g\_out};
  \item A última \textit{thread} a terminar posta \texttt{sem\_done};
  \item \texttt{img\_worker} grava \texttt{g\_out} em disco com \texttt{write\_pgm}.
\end{enumerate}
 
\section{Implementação}

\subsection{Processo Emissor (Sender)}

\textbf{Núcleo do emissor -- \texttt{img\_sender.c}}
\begin{lstlisting}[language=C, texcl=true]

// Garante existência da FIFO
    if(mkfifo(args.fifo, 0666) == -1 && errno != EEXIST){
    perror("Error creating FIFO.");
    return 1;
}
  
// Lê a imagem PGM do disco, decompõe seus metadados e pixels, e os separa nas variáveis PGM e Header
PGM pgm;
read_pgm(args.input_file, &pgm);

Header header;
header.w = pgm.w;
header.h = pgm.h;
header.maxv = pgm.maxv;
  
// Abre a FIFO para escrita em binário, bloqueando até que o processo trabalhador a abra para leitura
FILE* fd = fopen(args.fifo, "wb");
if(fd == NULL){
    perror("Error opening FIFO for writing.");
    free(pgm.data);
    return 1;
}

// Garante o envio dos metadados e os pixels da imagem para a FIFO
if(fwrite(&header, sizeof(Header), 1, fd) != 1){
    fprintf(stderr, "Error writing header to FIFO.\n");
    fclose(fd);
    free(pgm.data);
    return 1;
}

size_t img_size = ((size_t)pgm.w * (size_t)pgm.h);
if(fwrite(pgm.data, 1, img_size, fd) != img_size){
    fprintf(stderr, "Error writing image data to FIFO.\n");
    fclose(fd);
    free(pgm.data);
    return 1;
}

// Fecha FIFO e libera recursos alocados
free(pgm.data);
fflush(fd);
fclose(fd);
\end{lstlisting}

\subsection{Processo Trabalhador (Worker)}

\subsection{Filtros}
 
\begin{lstlisting}[language=C, texcl=true]
void apply_negative_block(int rs, int re){
  size_t i;
  for(size_t row = rs; row <= re; row++){
    for(size_t column = 0; column < g_in.w; column++){
      i = (g_in.w * row) + column;
      g_out.data[i] = (g_in.maxv - g_in.data[i]);
    }
  }
}

void apply_slice_block(int rs, int re, int t1, int t2){
  size_t i;
  for(size_t row = rs; row < re; row++){
    for(size_t column = 0; column < g_in.w; column++){
      i = (g_in.w * row) + column;
      if(g_in.data[i] >= t1 && g_in.data[i] <= t2){
        g_out.data[i] = g_in.data[i];
      }else{
        g_out.data[i] = 0;
      }
    }
  }
}
\end{lstlisting}

\section{Resultados}

\subsection{Correção dos Filtros}

Os filtros foram validados visualmente através do VIMIV, comparando a
imagem original às imagens processadas. O filtro negativo produziu a
inversão correta dos tons de cinza conforme a equação. O filtro
de limiarização com fatiamento preservou os pixels pertencentes ao intervalo
$[t_1, t_2]$ e zerou os demais, em conformidade com a equação.

\end{document}
